% Part: normal-modal-logic
% Chapter: frame-correspondence
% Section: first-order-definability

\documentclass[../../../include/open-logic-section]{subfiles}

\begin{document}

\olfileid{nml}{frd}{fol}

\olsection{प्रथम-क्रमातील परिभाष्यता}

चौकटींच्या प्राप्यता संबंधांचे अनेक गुणधर्म मोडल
!!{formula}s नी परिभाषित करता येतात, हे आपण पाहिले.
उदाहरणार्थ, चौकटींची सममिती \Ax{B}, म्हणजे
$p \lif \Box\Diamond
p$, या !!{formula} ने परिभाषित होते. आतापर्यंत
पाहिलेल्या सर्व अटी एकच द्विस्थानी !!{predicate}
असलेल्या !!a{language} तील प्रथम-क्रम !!{formula}s
नी व्यक्त करता येतात. उदाहरणार्थ, सममिती
$\lforall[x][\lforall[y][(\Atom{Q}{x,y} \lif \Atom{Q}{y, x})]]$
याने परिभाषित होते: $\Domain{M} =
W$ आणि $\Assign{Q}{M} = R$ असलेली प्रथम-क्रम
!!{structure}~$\Struct{M}$ वरील सूत्राची पूर्ति करते
तेव्हा आणि केवळ तेव्हाच, जेव्हा $R$ सममित असतो.
यातून पुढील व्याख्या सुचते:

\begin{defn}
  चौकटींचा वर्ग~$\mClass{F}$ \emph{प्रथम-क्रमात
  परिभाष्य} असतो, जर एकच द्विस्थानी !!{predicate}~$Q$
  असलेल्या प्रथम-क्रम भाषेतील असे !!a{sentence}~$!A$
  असेल की, $\Domain{M} = W$ आणि $\Assign{Q}{M}
  = R$ असलेल्या प्रथम-क्रम !!{structure}~$\Struct{M}$
  मध्ये $\Sat{M}{!A}$ तेव्हा आणि केवळ तेव्हाच,
  जेव्हा $\mModel{F} =
  \tuple{W, R} \in \mClass{F}$.
\end{defn}

आतापर्यंत पाहिलेले गुणधर्म आणि त्यांना परिभाषित करणारी
मोडल !!{formula}s ही अपवादात्मक उदाहरणे ठरतात.
प्रत्येक !!{formula} प्रथम-क्रमात परिभाष्य चौकटींचाच
वर्ग परिभाषित करते असे नाही; तसेच प्रथम-क्रमात
परिभाष्य चौकटींचा प्रत्येक वर्ग मोडल सूत्राने
परिभाषित होतो असेही नाही.

पहिल्या दाव्याचे प्रतिउदाहरण L\"ob सूत्र देते:
\begin{equation}
\Box(\Box p \lif p) \lif \Box p. \tag{\Ax{W}}
\end{equation}
\Ax{W} संक्रमक आणि प्रतिलोम-सुस्थापित चौकटींचा वर्ग
परिभाषित करते. $Rw_2w_1$, $Rw_3w_2$, \dots{}
अशी $w_1$, $w_2$, \dots{} ही अनंत श्रेणी नसल्यास
संबंध सुस्थापित असतो. उदाहरणार्थ, $\Nat$ वरील
$<$ हा संबंध सुस्थापित आहे; पण $\Int$ वरील
$<$ सुस्थापित नाही. एखाद्या संबंधाचा प्रतिलोम
सुस्थापित असेल, तर आणि तेव्हाच, तो संबंध प्रतिलोम-सुस्थापित
असतो. म्हणजे प्रतिलोम-सुस्थापित संबंधात $Rw_1w_2$,
$Rw_2w_3$, \dots{} अशी $w_1$, $w_2$, \dots{}
अनंत श्रेणी नसते.

मात्र संक्रमक प्रतिलोम-सुस्थापित संबंध परिभाषित करणारे
प्रथम-क्रम !!{formula} नाही. असे काही सूत्र असल्याचे
समजून, $R =
\Assign{Q}{M}$ संक्रमक आणि प्रतिलोम-सुस्थापित असेल
तेव्हा आणि केवळ तेव्हाच $\Sat{M}{!F}$, असे माना.
नवे स्थिरांक असलेल्या विस्तारित भाषेत $!A_n$ हे
पुढील !!{formula} असू द्या:
\[
(\Atom{Q}{a_1,a_2} \land \dots \land
    \Atom{Q}{a_{n},a_{n+1}})
\]
आता पुढील !!{formula}s चा संच पाहू:
\[
\Gamma = \{!F, !A_1, !A_2, \dots\}.
\]
$\Gamma$ चा प्रत्येक सांत उपसंच पूर्ततायोग्य आहे.
उपसंचात $!A_k$ असेल, तर $k$ हा त्यातील सर्वात मोठा
निर्देशांक घ्या; असे सूत्र नसेल, तर त्याचे मूल्य एक
घ्या. $\Domain{M_k} = \{1, \dots, k+1\}$,
आवश्यक स्थिरांकांसाठी $\Assign{a_i}{M_k} = i$,
आणि $\Assign{Q}{M_k} = <$ असे ठेवा; उर्वरित
स्थिरांकांना प्रांतातील कोणतेही मूल्य द्या.
$\{1,
\dots, k+1\}$ वरील $<$ संक्रमक आणि प्रतिलोम-सुस्थापित
असल्यामुळे $\Sat{M_k}{!F}$. रचनेनुसार प्रत्येक
$i \le k$ साठी $\Sat{M_k}{!A_i}$.
प्रथम-क्रम तर्कशास्त्राच्या संहतता प्रमेयावरून
$\Gamma$ ची कोणत्यातरी !!{structure}~$\Struct{M}$
मध्ये पूर्ति होते. गृहीतकावरून $\Sat{M}{!F}$
असल्यामुळे $\Assign{Q}{M}$ हा संबंध प्रतिलोम-सुस्थापित
आहे. पण $\Assign{a_1}{M}$, $\Assign{a_2}{M}$,
\dots{} यांपासून प्रतिलोम-सुस्थापितपणाने निषिद्ध
ठरवलेली अनंत श्रेणी मिळते.

दुसऱ्या दाव्याचे प्रतिउदाहरण म्हणजे वैश्विकता:
प्रत्येक $u$ आणि $v$ साठी $Ruv$. वैश्विक चौकटी
$\lforall[x][\lforall[y][\Atom{Q}{x,y}]]$ या
सूत्राने प्रथम-क्रमात परिभाषित होतात. मात्र एखादे
मोडल सूत्र नेमक्या वैश्विक चौकटींमध्येच वैध
असते असे नाही. हा निष्कर्ष स्वतंत्रपणे महत्त्वाच्या
एका तथ्यावरून येतो: वैश्विक चौकटींमध्ये वैध
असलेली सूत्रे आणि परावर्ती, सममित व संक्रमक
चौकटींमध्ये वैध असलेली सूत्रे नेमकी तीच आहेत.
परावर्ती, सममित आणि संक्रमक असूनही वैश्विक
नसलेल्या चौकटी आहेत. म्हणून सर्व वैश्विक चौकटींमध्ये
वैध असलेले प्रत्येक !!{formula} काही अवैश्विक
चौकटींमध्येही वैध असते.

\end{document}
